\documentclass[10pt]{article}
\usepackage[letterpaper,margin=0.74in]{geometry}
\usepackage{amsmath,amssymb,amsthm,booktabs,microtype}
\usepackage[T1]{fontenc}
\usepackage{lmodern}
\usepackage[hidelinks]{hyperref}

\newtheorem{theorem}{Theorem}
\newtheorem{proposition}[theorem]{Proposition}
\newtheorem{lemma}[theorem]{Lemma}
\newtheorem{corollary}[theorem]{Corollary}
\newtheorem{remark}[theorem]{Remark}
\newcommand{\Z}{\mathbb Z}
\newcommand{\C}{\mathbb C}
\newcommand{\one}{\mathbf 1}
\newcommand{\ip}[2]{\langle #1,#2\rangle}
\newcommand{\Ax}{A_x}
\newcommand{\Qx}{\mathcal Q_x}
\newcommand{\Li}{\operatorname{Li}}
\newcommand{\TV}{\operatorname{TV}}
\newcommand{\e}{\varepsilon}
\newcommand{\zzero}{z_0}
\newcommand{\zone}{z_1}
\newcommand{\ztwo}{z_2}
\setlength{\emergencystretch}{3em}

\title{Four-Block Haar Lifts and a Transverse Norm Floor\
for the Literal V59 Adjoint}
\author{Liang Wang\\
\small School of Mathematics and Statistics, Huazhong University of Science and Technology (HUST),\\
\small Wuhan, China}
\date{August 26, 2026}

\begin{document}
\maketitle

\begin{abstract}
We test whether the nonzero ordered-rank Haar coefficient found in TPC-256
could leave the orthogonal output of the literal V59 adjoint negligible.  We
split each of the two rank children into two consecutive blocks, before
inspecting any coefficient, and form the global midpoint contrast together
with two within-child Haar contrasts.  These three vectors are exactly
orthonormal at every admissible finite clock.  Second-order prime-density
curvature gives three explicit positive beta moments, while the exact
TPC-255 diagonal/boundary compiler gives the same-order asymptotic
\[
 \ip{z_i}{A_x\beta}
 =-\left(\frac92\kappa_i+o(1)\right)
   \frac{x^{7/6}}{\log^3 x}, \qquad i=0,1,2.
\]
Here \(\kappa_0=\log(32/27)/\sqrt2\),
\(\kappa_1=\log(3456/3125)/2\), and
\(\kappa_2=\log(884736/823543)/2\).  Parseval therefore supplies an
explicit same-order lower floor in the two-dimensional plane orthogonal to
the old midpoint.  This is a transverse obstruction, not an upper \(L^2\)
estimate: full output control, full Gate B, the strict global \(1/400\)
payment, and a twin-prime conclusion remain open.
\end{abstract}

\section{Question and claim boundary}

The preceding rank-midpoint computation isolated a genuine literal scalar
\cite{tpc256},
but one scalar does not determine a vector.  The precise question here is
whether the orthogonal component can nevertheless be lower order for the same
operator and coefficient.  We answer the smallest source-only version of
that question: two descendants of the left and right rank children.

The word ``floor'' is intentional.  We prove that a finite projection is at
least of a stated size.  We do not prove that the unprojected vector is at
most that size.  In particular, no arithmetic \(L^2\) upper bound is hidden
in the result.

\section{Literal clock and four-block frame}

Let \(x\to\infty\) through real values and put
\begin{equation}
 a=\lfloor x/2\rfloor,\qquad b=\lfloor x\rfloor,
 \qquad I_x=\{a+1,\ldots,b\},\qquad N=b-a.
 \label{eq:clock}
\end{equation}
The ordered rank split is inherited from the source-frozen compiler
\cite{tpc253}:
\begin{equation}
 \ell=\lfloor N/2\rfloor,\quad r=N-\ell,\quad m=a+\ell,
 \quad L=\{a+1,\ldots,m\},\quad R=\{m+1,\ldots,b\}.
 \label{eq:rank}
\end{equation}
Split the two children in their physical order:
\begin{equation}
 s_1=\lfloor\ell/2\rfloor,\quad s_2=\ell-s_1,
 \qquad s_3=\lfloor r/2\rfloor,\quad s_4=r-s_3.       \label{eq:blocksizes}
\end{equation}
Let \(B_1,B_2\) be the consecutive subintervals of \(L\) of lengths
\(s_1,s_2\), and let \(B_3,B_4\) be the analogous subintervals of \(R\).
For adjacent intervals \(A,B\) of lengths \(p,q\), define
\begin{equation}
 h(A,B)=\left(\frac{pq}{p+q}\right)^{1/2}
 \left(\frac{\one_A}{p}-\frac{\one_B}{q}\right).       \label{eq:contrast}
\end{equation}
Our frame is
\begin{equation}
 \zzero=h(B_1\cup B_2,B_3\cup B_4),\qquad
 \zone=h(B_1,B_2),\qquad \ztwo=h(B_3,B_4).              \label{eq:frame}
\end{equation}
The vector \(\zzero\) is exactly the TPC-256 midpoint vector.

\begin{lemma}[Exact orthonormality and variation]\label{lem:frame}
For all sufficiently large \(x\), the three vectors in \eqref{eq:frame}
are source-only, exactly orthonormal, and, after zero extension outside
\(I_x\), satisfy
\[
 \TV(z_i)=\frac{2}{\rho_i},
 \qquad \rho_i^2=\frac{p_iq_i}{p_i+q_i},
\]
where \((p_0,q_0)=(\ell,r)\), \((p_1,q_1)=(s_1,s_2)\), and
\((p_2,q_2)=(s_3,s_4)\).
\end{lemma}

\begin{proof}
For a pair of sizes \(p,q\), direct counting gives
\[
 \|h(A,B)\|_2^2=\frac{pq}{p+q}\left(\frac1p+\frac1q\right)=1.
\]
The sum of \(\zone\) over \(B_1\cup B_2\) is zero, whereas \(\zzero\) is
constant there.  Thus \(\ip{\zzero}{\zone}=0\).  The same argument on the
right gives \(\ip{\zzero}{\ztwo}=0\), and the last two vectors have disjoint
support.  This proves orthonormality, including all floor-rounding cases.

For the zero extension of a contrast, the successive nonzero levels are
\(\rho_i/p_i\) and \(-\rho_i/q_i\).  The two outer jumps and the internal
jump have total size
\[
 \frac{\rho_i}{p_i}+\rho_i\left(\frac1{p_i}+\frac1{q_i}\right)
       +\frac{\rho_i}{q_i}
 =\frac{2}{\rho_i}.
\]
The construction uses only ordered coordinates and the clock, so it is
coefficient-independent by definition.
\end{proof}

The scaled block limits are
\[
 B_1=[1/2,5/8],\ B_2=[5/8,3/4],\quad
 B_3=[3/4,7/8],\ B_4=[7/8,1],
\]
up to endpoint errors of size \(O(1/x)\).  In particular all three
normalizations are comparable with \(\sqrt{x}\).

\section{The literal coefficient and its three curvatures}

Retain the V59 scales and literal coefficient
\begin{equation}
 H=x^{21/32},\qquad Q=x^{1/3},\qquad U=x^{133/400},
 \qquad \beta(t)=\frac{\Lambda(t)}{\log t}
 -\sum_{\substack{d\mid t\\d\leq U}}\mu(d).             \label{eq:beta}
\end{equation}
For an interval \(J\) of length \(s\),
\[
 \#\{n\in J:d\mid n\}=s/d+\theta_{J,d},\qquad
 |\theta_{J,d}|\leq 1.
\]
The leading \(1/d\) term cancels for every divisor between the two children
of each contrast.  Consequently
\begin{equation}
 \left|\ip{z_i}{\sum_{d\mid\cdot,\ d\leq U}\mu(d)}\right|
 \leq \frac{U}{\rho_i}=O(x^{-67/400}).                 \label{eq:divisor}
\end{equation}
No Mertens cancellation is used.

Let \(F(y)=\sum_{2\leq n\leq y}\Lambda(n)/\log n\).  Prime-power
decomposition and the source-locked de la Vallée Poussin theorem
\cite{davenport,montgomery-vaughan,ik,tpc233} give
\begin{equation}
 F(y)=\Li(y)+O\left(y\exp(-c\sqrt{\log y})\right).     \label{eq:pnt}
\end{equation}
For equal-width limiting intervals \(A,B\), of width \(d\), expanding
\(1/\log(xy)\) gives
\begin{equation}
 \operatorname{mean}_A(P)-\operatorname{mean}_B(P)
 =\frac{1}{d\log^2x}
   \left(\int_B\log y\,dy-\int_A\log y\,dy\right)
   +O(\log^{-3}x),                                    \label{eq:curvature}
\end{equation}
where \(P=\Lambda/\log\).  The endpoint and PNT errors are absorbed in the
remainder.  The three elementary integrals are summarized in Table~\ref{tab:curv}.

\begin{table}[ht]
\centering
\caption{Second-order curvature constants.}
\label{tab:curv}
\begin{tabular}{@{}ccccc@{}}
\toprule
vector & limiting pair & mean curvature & $\rho_i/\sqrt{x}$ & $\kappa_i$\\
\midrule
$z_0$ & $[1/2,3/4],[3/4,1]$ & $2\log(32/27)$ & $1/(2\sqrt2)$ & $\log(32/27)/\sqrt2$\\
$z_1$ & $[1/2,5/8],[5/8,3/4]$ & $2\log(3456/3125)$ & $1/4$ & $\log(3456/3125)/2$\\
$z_2$ & $[3/4,7/8],[7/8,1]$ & $2\log(884736/823543)$ & $1/4$ & $\log(884736/823543)/2$\\
\bottomrule
\end{tabular}
\end{table}

All three logarithm arguments exceed one.  Combining the table with
\eqref{eq:divisor} proves
\begin{equation}
 \ip{z_i}{\beta}=(\kappa_i+O(1/\log x))\frac{\sqrt{x}}{\log^2x}>0
 \quad (i=0,1,2)                                  \label{eq:betacontrasts}
\end{equation}
for all sufficiently large real \(x\).

\section{Adjoint normal form for bounded variation}

Let \(\Qx=\{q\text{ prime}:Q<q\leq2Q\}\), and let
\(K_H(h)=\widehat{\psi_+}(h/H)\), with the fixed smooth compactly supported
profile normalized by \(\int\psi_+=1\).  The literal operator is
\begin{equation}
 \Ax(u,t)=\one_{u\ne t}\sum_{q\in\Qx}q\one_{q\nmid u}\one_{q\nmid t}
 K_H(u-t)\left(\one_{u\equiv t\ (q)}-\frac1{q-1}\right). \label{eq:operator}
\end{equation}
For \(q\nmid t\), retain the combined row from the exact adjoint compiler
\cite{tpc255}:
\[
 v_{q,t}(u)=\one_{q\nmid u}\left(\one_{u\equiv t\ (q)}-\frac1{q-1}\right).
\]
Its complete period is zero only after the two displayed pieces are
 recombined.  The exact TPC-255 Poisson calculation, based on the V43
transference \cite{v43}, applies to every
zero-extended piecewise-constant test function and yields
\begin{equation}
 \ip{z}{\Ax\beta}=-B_Q\ip{z}{\beta}+R_{\rm unit}(z)+R_{\rm bdry}(z),
 \qquad B_Q=\sum_{q\in\Qx}\frac{q(q-2)}{q-1}.             \label{eq:normalform}
\end{equation}
The term \(R_{\rm bdry}\) includes every outer endpoint and every internal
jump of the zero extension.  This is an exact algebraic extension of the
one-jump identity; it does not discard a boundary.

The combined row obeys
\begin{equation}
 |v_{q,t}(t+h)|\leq\one_{q\mid h}+\frac2q,\qquad
 \sum_h|hK_H(h)|\left(\one_{q\mid h}+\frac2q\right)
 \ll_{\psi_+}\frac{H^2}{q}.                           \label{eq:firstmoment}
\end{equation}
For a fixed displacement \(h\), at most \(|h|\) source points cross any
one jump.  Thus, for a fixed \(\e>0\),
\begin{equation}
 |R_{\rm bdry}(z)|\ll_{\psi_+,\e}QH^2\TV(z)x^\e,
 \qquad |R_{\rm unit}(z)|\ll_\e x^{5/6+\e}             \label{eq:remainders}
\end{equation}
for each of the three normalized contrasts.  Since \(\TV(z_i)=O(x^{-1/2})\),
the first quantity is \(O_{\psi_+,\e}(x^{55/48+\e})\).

\begin{lemma}[Diagonal scale]\label{lem:bq}
The weighted prime-shell input gives
\[
 B_Q=\left(\frac92+o(1)\right)\frac{x^{2/3}}{\log x}.
\]
\end{lemma}

\begin{proof}
Use \(q(q-2)/(q-1)=q-1-1/(q-1)\) and the weighted prime number theorem
recorded in the top-shell ledger \cite{v59}:
\[
 \sum_{Q<q\leq2Q}q=\left(\frac32+o(1)\right)\frac{Q^2}{\log Q}.
\]
The lower-order terms are absorbed, and \(Q=x^{1/3}\) gives the claim.
\end{proof}

\section{Main theorem and norm floors}

\begin{theorem}[Three-mode adjoint asymptotics]\label{thm:main}
For each \(i=0,1,2\),
\begin{equation}
 \boxed{\displaystyle
 \ip{z_i}{\Ax\beta}
 =-\left(\frac92\kappa_i+o(1)\right)
   \frac{x^{7/6}}{\log^3x}}\qquad\text{in }\C.              \label{eq:scalar}
\end{equation}
\end{theorem}

\begin{proof}
Insert \eqref{eq:betacontrasts} and Lemma~\ref{lem:bq} into
\eqref{eq:normalform}.  The diagonal exponent is
\(2/3+1/2=7/6=56/48\).  The boundary exponent from \eqref{eq:remainders} is
\[
 1/3+2(21/32)-1/2=55/48,
\]
and the input-unit exponent is \(5/6\).  Choose a fixed
\(0<\e<1/48\).  All remainders are then lower order than the diagonal, and
the asserted complex asymptotic follows.  The proof uses only absolute
remainder bounds after the exact combined-row cancellation; it assumes no
reality or evenness of the kernel.
\end{proof}

\begin{corollary}[Transverse and three-mode floors]\label{cor:floors}
Let \(Z=\operatorname{span}\{z_0,z_1,z_2\}\) and
\(T=\operatorname{span}\{z_1,z_2\}\).  Then
\begin{align}
 \|P_Z\Ax\beta\|_2
 &=\left(\frac92\sqrt{\kappa_0^2+\kappa_1^2+\kappa_2^2}+o(1)\right)
   \frac{x^{7/6}}{\log^3x},                              \label{eq:zfloor}\\
 \|P_T\Ax\beta\|_2
 &=\left(\frac92\sqrt{\kappa_1^2+\kappa_2^2}+o(1)\right)
   \frac{x^{7/6}}{\log^3x}.                              \label{eq:tfloor}
\end{align}
In particular, \(T\subset z_0^\perp\), so the second line is a lower floor
for \(\|(I-z_0\otimes z_0)\Ax\beta\|_2\).
\end{corollary}

\begin{proof}
Lemma~\ref{lem:frame} gives an orthonormal family.  Therefore finite
dimensional Parseval gives
\[
 \|P_Zg\|_2^2=\sum_{i=0}^2|\ip{z_i}{g}|^2,
 \qquad \|P_Tg\|_2^2=\sum_{i=1}^2|\ip{z_i}{g}|^2.
\]
Apply this with \(g=\Ax\beta\) and use Theorem~\ref{thm:main}.  The
constants are positive, so taking square roots is stable.  The inclusion of
\(T\) in \(z_0^\perp\) was proved exactly in Lemma~\ref{lem:frame}.
\end{proof}

The factors before \(9/2\) are approximately
\[
 0.135096662713318\quad\text{and}\quad0.061792126717520,
\]
respectively.  These decimals are orientation data only; the logarithmic
forms in the theorem are the actual constants.

\section{Interpretation, obstruction, and next route}

The result establishes a source-backed same-order transverse signal for the
literal object.  It therefore refutes the shortcut
\[
 \text{one nonzero midpoint coefficient}
 \Longrightarrow
 \text{the entire transverse output is lower order}.
\]
The refutation is scoped: it concerns this literal V59 clock and this
coefficient-independent four-block frame, not arbitrary operators.

Several gates remain untouched.  The theorem gives no upper bound for
\(\|\Ax\beta\|_2\), no uniform estimate over a growing Haar basis, no signed
coupling to the physical \(w\) lane, and no strict global \(1/400\) payment.
The correct next question is whether the two transverse leading constants
can be cancelled by a source-frozen linear combination.  That null-direction
test may reveal a logarithmic or boundary-scale remainder, but it must be
proved separately.

\section*{Reproducibility and epistemic labels}

The companion certificate checks the finite rational identities, source hashes,
curvature logarithm vectors, and exponent ledger.  Its finite beta samples are
marked \texttt{NUMERICAL\_OBSERVATION}; they are not evidence for the PNT
asymptotic.  The status of this paper is
\[
\texttt{PROVED\_SOURCE\_BACKED\_TRANSVERSE\_HAAR\_NORM\_FLOOR\_FOR\_LITERAL\_V59\_ADJOINT}.
\]

\bibliographystyle{plain}
\bibliography{references}
\end{document}
